%==============================================================================
% PINN-ONB01 — Supplementary Material
% Companion to the main manuscript on
% "Surface-Conditioned Physics-Informed Neural Network for Pool-Boiling
%  Onset of Nucleate Boiling"
%
% Build with:
%   latexmk -pdf supplementary.tex
%==============================================================================
\documentclass[11pt, a4paper]{article}

\usepackage[utf8]{inputenc}
\usepackage[T1]{fontenc}
\usepackage{lmodern}
\usepackage{amsmath, amssymb, amsfonts}
\usepackage[separate-uncertainty=true]{siunitx}
% Match the main manuscript: comma thousands separator for 4-digit+ numerals.
\sisetup{group-separator={,}, group-minimum-digits=4, group-digits=integer, output-decimal-marker={.}}
\usepackage{graphicx}
\graphicspath{{../figures/}{../../04_analysis/figures/}}
\usepackage{booktabs}
\usepackage[margin=2.5cm]{geometry}
\usepackage{hyperref}
\hypersetup{colorlinks=true, linkcolor=blue, citecolor=blue, urlcolor=blue}
\usepackage[capitalise, noabbrev]{cleveref}

\title{Supplementary Material\\
\large Surface-Conditioned Physics-Informed Neural Network for
Pool-Boiling Onset of Nucleate Boiling}
\author{}
\date{}

\renewcommand{\thefigure}{S\arabic{figure}}
\renewcommand{\thetable}{S\arabic{table}}
\renewcommand{\theequation}{S\arabic{equation}}

\begin{document}
\maketitle

\section*{Overview}

This supplement collects per-test figures that complement the composite
panels in the main manuscript. \cref{sec:supp_level1} unpacks the
Level~1 verification composite into its four constituent tests, and
\cref{sec:supp_physics} unpacks the physics-consistency composite into
its five trend sweeps. \cref{sec:supp_inverse} shows the PINN-augmented
inverse parity plot, which the main text references but does not
display. \cref{sec:supp_dataset} reports extended dataset statistics
omitted from the main text.

% =============================================================================
\section{Level 1 verification (single-test figures)}\label{sec:supp_level1}
% =============================================================================

The main manuscript condenses the four canonical numerical-verification
tests of the PINN backbone into a single composite figure. For
completeness, the original single-test figures are reproduced below.

\begin{figure}[ht]
  \centering
  \includegraphics[width=0.65\linewidth]{level1_1d_conduction_parity}
  \caption{\textbf{Level~1, Test~1.} Parity between the PINN
    prediction and the analytical 1D steady heat-conduction profile;
    relative $L^{2}$ error \SI{0.018}{\percent}.}
  \label{fig:supp_level1_conduction}
\end{figure}

\begin{figure}[ht]
  \centering
  \includegraphics[width=0.65\linewidth]{level1_pde_residual_convergence}
  \caption{\textbf{Level~1, Test~2.} Monotonic convergence of the PDE
    residual loss as the number of collocation points increases
    ($N_{\mathrm{coll}}\in\{200,\ldots,5\,000\}$; log--log slope
    $+0.009$).}
  \label{fig:supp_level1_pde}
\end{figure}

\begin{figure}[ht]
  \centering
  \includegraphics[width=0.65\linewidth]{level1_nc_nu_comparison}
  \caption{\textbf{Level~1, Test~3.} PINN natural-convection prediction
    against the McAdams correlation
    $\mathrm{Nu} = 0.54\,\mathrm{Ra}^{1/4}$; relative error
    \SI{0.000}{\percent}.}
  \label{fig:supp_level1_nc}
\end{figure}

\begin{figure}[ht]
  \centering
  \includegraphics[width=0.65\linewidth]{level1_autograd_accuracy}
  \caption{\textbf{Level~1, Test~4.} Autograd-based first and second
    derivatives of a manufactured $\sin(2\pi z)$ field, with
    \SI{0.026}{\percent} and \SI{0.011}{\percent} relative error,
    respectively.}
  \label{fig:supp_level1_autograd}
\end{figure}

% =============================================================================
\section{Physics-consistency trends (single-variable sweeps)}\label{sec:supp_physics}
% =============================================================================

The five physical-trend sweeps that make up the audit figure in the
main manuscript are reproduced here as individual panels. Each panel
plots the \texttt{baseline\_phaseDbal} ensemble mean of
$\Delta T_{\mathrm{ONB}}$ against a single varied descriptor, with the
remaining descriptors held at the dataset median.

\begin{figure}[ht]
  \centering
  \includegraphics[width=0.62\linewidth]{physics_trend_qflux}
  \caption{\textbf{Physics trend, heat flux.} Synthetic heat-flux sweep:
    $q''\uparrow \Rightarrow \Delta T_{\mathrm{ONB}}\uparrow$
    (Hsu $\sqrt{q''}$ trend); Spearman $\rho = +0.970$.}
  \label{fig:supp_physics_qflux}
\end{figure}

\begin{figure}[ht]
  \centering
  \includegraphics[width=0.62\linewidth]{physics_trend_roughness}
  \caption{\textbf{Physics trend, surface roughness.} Synthetic roughness sweep:
    $R_{a}\uparrow \Rightarrow \Delta T_{\mathrm{ONB}}\downarrow$;
    Spearman $\rho = -0.815$.}
  \label{fig:supp_physics_roughness}
\end{figure}

\begin{figure}[ht]
  \centering
  \includegraphics[width=0.62\linewidth]{physics_trend_contact_angle}
  \caption{\textbf{Physics trend, contact angle.} Synthetic contact-angle sweep:
    $\theta\uparrow \Rightarrow \Delta T_{\mathrm{ONB}}\downarrow$;
    Spearman $\rho = -0.720$.}
  \label{fig:supp_physics_theta}
\end{figure}

\begin{figure}[ht]
  \centering
  \includegraphics[width=0.62\linewidth]{physics_trend_pressure}
  \caption{\textbf{Physics trend, pressure.} Synthetic pressure sweep:
    $P\uparrow \Rightarrow \Delta T_{\mathrm{ONB}}\downarrow$;
    Spearman $\rho = -1.000$.}
  \label{fig:supp_physics_pressure}
\end{figure}

\begin{figure}[ht]
  \centering
  \includegraphics[width=0.62\linewidth]{physics_trend_subcool}
  \caption{\textbf{Physics trend, subcooling.} Synthetic subcooling sweep:
    $\Delta T_{\mathrm{sub}}\uparrow \Rightarrow
    \Delta T_{\mathrm{ONB}}\uparrow$;
    Spearman $\rho = +1.000$.}
  \label{fig:supp_physics_subcool}
\end{figure}

% =============================================================================
\section{PINN-augmented inverse parity}\label{sec:supp_inverse}
% =============================================================================

The main text summarizes the PINN-augmented inverse cavity-radius
result. The corresponding parity plot is provided below.

\begin{figure}[ht]
  \centering
  \includegraphics[width=0.65\linewidth]{inverse_pinn_vs_hsu}
  \caption{\textbf{PINN-augmented inverse parity.} Estimated $r_{c}$
    against the Hsu-analytical inverse on the converged subset ($n=18$);
    Spearman $\rho = +0.34$, $p = 0.17$. The remaining 30 surfaces
    collapsed onto the initial guess and are omitted.}
  \label{fig:supp_inverse_pinn_hsu}
\end{figure}

% =============================================================================
\section{Extended dataset statistics}\label{sec:supp_dataset}
% =============================================================================

\Cref{tab:supp_surface_inventory} lists the surface-card categories
used in the dataset, together with the number of distinct surfaces and
ONB observations per category. The script
\verb|02_data/scripts/update_surface_id.py| maintains the mapping from
raw figure series to surface-card identifier, and the full mapping is
released as part of the open dataset.

\begin{table}[ht]
\centering
\caption{Surface-card categories used in the dataset.}
\label{tab:supp_surface_inventory}
\small
\begin{tabular}{lcc}
\toprule
Category & Surfaces & ONB labels \\
\midrule
\texttt{betz}        &  7 & 10 \\
\texttt{bourdon12}   &  5 &  6 \\
\texttt{bourdon15}   &  4 &  5 \\
\texttt{jabardo}     &  8 & 10 \\
\texttt{jabardo\_br} &  5 & 10 \\
\texttt{jabardo\_ss} &  7 & 14 \\
\texttt{jones}       &  5 & 10 \\
\texttt{jones\_w}    &  1 &  0 \\
\texttt{jo}          &  2 &  2 \\
\texttt{phan}        &  5 &  5 \\
\midrule
Total                & 49 & 82 \\
\bottomrule
\end{tabular}
\end{table}

% =============================================================================
\section{Hyperparameter optimization}\label{sec:supp_hpo}
% =============================================================================

The loss weights, learning rate, and architectural integers reported in
the main manuscript were selected through a Tree-structured Parzen
Estimator (TPE) study run with the Optuna framework. The study explored
\num{30} trials, of which \num{21} completed successfully and \num{9}
were pruned for early-stopping criteria. The total wall-clock budget
was approximately \num{3.1} hours on a single Apple~M1 CPU.
\Cref{tab:supp_hpo_top5} lists the five best trials ranked by
validation RMSE on the \num{77}-point held-out ONB subset. The chosen
configuration (\texttt{baseline\_phaseDbal}) inherits the median values
of the top tier and refines the monotonicity-loss balance separately.

\begin{table}[ht]
\centering
\caption{Top five Optuna TPE trials ranked by validation RMSE on
  the held-out ONB subset.}
\label{tab:supp_hpo_top5}
\small
\begin{tabular}{cccccccc}
\toprule
Rank & Trial & val.\ RMSE [K] & test RMSE [K] &
$w_{\mathrm{cond}}$ & $w_{\mathrm{BC}}$ & $w_{\mathrm{data}}$ & $w_{\mathrm{ONB}}$ \\
\midrule
1 & 21 & 5.505 & 7.205 & 0.227 & 3.892 & 35.57 & 0.1014 \\
2 & 10 & 5.507 & 7.204 & 0.258 & 3.083 & 38.84 & 0.1071 \\
3 & 13 & 5.508 & 7.204 & 0.805 & 2.890 & 47.37 & 0.1002 \\
4 & 11 & 5.508 & 7.204 & 0.252 & 3.586 & 30.75 & 0.1026 \\
5 & 19 & 5.511 & 7.204 & 0.581 & 5.113 & 18.84 & 0.4285 \\
\bottomrule
\end{tabular}

\vspace{0.5em}
All five winners select \texttt{conditioning=film}, $d_z=8$,
hidden dim $64$, and three hidden layers, indicating that the
architectural choice is decisive while the loss weights remain in
a comparable basin.
\end{table}

The Optuna importance scores ranked the explored parameters as
hidden\_dim (0.30), $\eta_{\mathrm{Adam}}$ (0.22), $w_{\mathrm{data}}$
(0.21), $d_z$ (0.17), and $w_{\mathrm{ONB}}$ (0.07). The \num{2.0}
value of $w_{\mathrm{cond}}$ adopted in the final model is a rounded
upper bound that keeps the PDE residual within the same order of
magnitude as the data term during the warm-up stage.

% =============================================================================
\section{Model design progression}\label{sec:supp_phases}
% =============================================================================

The headline \texttt{baseline\_phaseDbal} model is the outcome of a
sequence of design iterations that refined the conditioning,
regularization, and loss balance. \Cref{tab:supp_phases} summarizes
the principal change between consecutive phases and the corresponding
validation RMSE on the test split. Phases C and E illustrate the
sensitivity of the model to over-regularization, while phase D and the
balanced variant Dbal converged on the published configuration.

\begin{table}[ht]
\centering
\caption{Model design progression. Each phase modifies the
  preceding configuration by the indicated change. Test RMSE values
  are end-of-Phase-3 validation metrics on the \num{77}-point held-out
  subset (last 200 epochs averaged).}
\label{tab:supp_phases}
\small
\begin{tabular}{lll}
\toprule
Phase & Principal change                       & Test RMSE [K] \\
\midrule
A     & Surface-aware PINN baseline (concat)   & 6.20 \\
B     & Switch to FiLM conditioning, HPO best  & 5.50 \\
C     & Add monotonicity penalty on $R_a$ only & 8.61 \\
C-A   & Reduce $w_{\mathrm{mono},R_a}$         & 7.20 \\
C-D   & Add monotonicity penalty on $\theta$   & 7.15 \\
D     & Rebalance data/PDE ratio               & 7.04 \\
Dbal  & Lower $w_{\mathrm{mono}}$, tightened
        scheduling (\textbf{published})        & \textbf{3.42} \\
E     & Aggressive $w_{\mathrm{ONB}}$ ramp     & 9.50 \\
\bottomrule
\end{tabular}

\vspace{0.5em}
Values for phases A through E are end-of-Phase-3 averages from
training logs (\texttt{03\_model/logs/phase\#\_p3\_*.log}).
The Dbal headline of \SI{3.42}{K} reflects the test-set evaluation
reported in Table~2 of the main text.
\end{table}

\end{document}
